%% kapa: exam=1 round=35 year=2024 period=2 form=A subject=appraisal_law subject_ko=감정평가관계법규 seq=5 number=5
%% source: Q-Net 감정평가사 기출문제 (cst00309) artlSeq=5213229 file=raw/1008/5213229/2교시 문제지 원본.pdf
%% conversion: kordoc 4.24.51 -> tools/gen1.py (auto); review=proofread-sonnet; flags=none
\begin{question}{5}
국토의 계획 및 이용에 관한 법령상 공동구에 관한 설명으로 옳은 것은?
\begin{choices}
  \item 「도시개발법」에 따른 100만제곱미터 규모의 도시개발구역에서 개발사업을 시행하는 자는 공동구를 설치하여야 한다.
  \item 통신선로는 공동구협의회의 심의를 거쳐야 수용할 수 있다.
  \item 공동구의 설치비용은 「국토의 계획 및 이용에 관한 법률」 또는 다른 법률에 특별한 규정이 있는 경우를 제외하고는 공동구 점용예정자와 사업시행자가 부담한다.
  \item 부담금의 납부통지를 받은 공동구 점용예정자는 공동구설치공사가 착수되기 전에 부담액의 3분의 2 이상을 납부하여야 한다.
  \item 공동구관리자는 1년에 2회 이상 공동구의 안전점검을 실시하여야 한다.
\end{choices}
\end{question}
\answer{3}
